\documentclass[10pt,a4paper]{amsart}
\usepackage{amsmath,amssymb,amsthm,graphicx}
\usepackage{verbatim}
\newenvironment{reference}{\comment}{\endcomment}
%\newenvironment{reference}{}{}
\newenvironment{slogan}{\comment}{\endcomment}
\newenvironment{history}{\comment}{\endcomment}
% For commutative diagrams we use Xy-pic
\usepackage[all]{xy}

% We use 2cell for 2-commutative diagrams.
\xyoption{2cell}
% Theorem environments.
%
\theoremstyle{plain}
\newtheorem{theorem}[subsection]{Theorem}
\newtheorem{proposition}[subsection]{Proposition}
\newtheorem{lemma}[subsection]{Lemma}

\theoremstyle{definition}
\newtheorem{definition}[subsection]{Definition}
\newtheorem{example}[subsection]{Example}
\newtheorem{exercise}[subsection]{Exercise}
\newtheorem{situation}[subsection]{Situation}

\theoremstyle{remark}
\newtheorem{remark}[subsection]{Remark}
\newtheorem{remarks}[subsection]{Remarks}

\numberwithin{equation}{subsection}

% Macros
%
\def\lim{\mathop{\mathrm{lim}}\nolimits}
\def\colim{\mathop{\mathrm{colim}}\nolimits}
\def\Spec{\mathop{\mathrm{Spec}}}
\def\Hom{\mathop{\mathrm{Hom}}\nolimits}
\def\Ext{\mathop{\mathrm{Ext}}\nolimits}
\def\SheafHom{\mathop{\mathcal{H}\!\mathit{om}}\nolimits}
\def\SheafExt{\mathop{\mathcal{E}\!\mathit{xt}}\nolimits}
\def\Sch{\mathit{Sch}}
\def\Mor{\mathop{\mathrm{Mor}}\nolimits}
\def\Ob{\mathop{\mathrm{Ob}}\nolimits}
\def\Sh{\mathop{\mathit{Sh}}\nolimits}
\def\NL{\mathop{N\!L}\nolimits}
\def\CH{\mathop{\mathrm{CH}}\nolimits}
\def\proetale{{pro\text{-}\acute{e}tale}}
\def\etale{{\acute{e}tale}}
\def\QCoh{\mathit{QCoh}}
\def\Ker{\mathop{\mathrm{Ker}}}
\def\Im{\mathop{\mathrm{Im}}}
\def\Coker{\mathop{\mathrm{Coker}}}
\def\Coim{\mathop{\mathrm{Coim}}}

% Boxtimes
%
\DeclareMathSymbol{\boxtimes}{\mathbin}{AMSa}{"02}

%
% Macros for moduli stacks/spaces
%
\def\QCohstack{\mathcal{QC}\!\mathit{oh}}
\def\Cohstack{\mathcal{C}\!\mathit{oh}}
\def\Spacesstack{\mathcal{S}\!\mathit{paces}}
\def\Quotfunctor{\mathrm{Quot}}
\def\Hilbfunctor{\mathrm{Hilb}}
\def\Curvesstack{\mathcal{C}\!\mathit{urves}}
\def\Polarizedstack{\mathcal{P}\!\mathit{olarized}}
\def\Complexesstack{\mathcal{C}\!\mathit{omplexes}}
% \Pic is the operator that assigns to X its picard group, usage \Pic(X)
% \Picardstack_{X/B} denotes the Picard stack of X over B
% \Picardfunctor_{X/B} denotes the Picard functor of X over B
\def\Pic{\mathop{\mathrm{Pic}}\nolimits}
\def\Picardstack{\mathcal{P}\!\mathit{ic}}
\def\Picardfunctor{\mathrm{Pic}}
\def\Deformationcategory{\mathcal{D}\!\mathit{ef}}
\usepackage{cleveref}
\crefname{originaldefinition}{original definition}{original definitions}
\Crefname{originaldefinition}{Original definition}{Original definitions}
\makeatletter
\newlabel{definition-fill-in-diagram}{{definition-fill-in-diagram}{source}{}{sourcepointer-1}{}}
\newlabel{definition-fill-in-diagram@cref}{{[originaldefinition][0][]definition-fill-in-diagram}{[1][0][]source}{}{}{}}
\makeatother
\crefname{originallemma}{original lemma}{original lemmas}
\Crefname{originallemma}{Original lemma}{Original lemmas}
\makeatletter
\newlabel{lemma-cat-dotted-arrows}{{lemma-cat-dotted-arrows}{source}{}{sourcepointer-2}{}}
\newlabel{lemma-cat-dotted-arrows@cref}{{[originallemma][0][]lemma-cat-dotted-arrows}{[1][0][]source}{}{}{}}
\makeatother
\crefname{originallemma}{original lemma}{original lemmas}
\Crefname{originallemma}{Original lemma}{Original lemmas}
\makeatletter
\newlabel{lemma-existence-representable}{{lemma-existence-representable}{source}{}{sourcepointer-3}{}}
\newlabel{lemma-existence-representable@cref}{{[originallemma][0][]lemma-existence-representable}{[1][0][]source}{}{}{}}
\makeatother
\crefname{originallemma}{original lemma}{original lemmas}
\Crefname{originallemma}{Original lemma}{Original lemmas}
\makeatletter
\newlabel{lemma-properties-diagonal}{{lemma-properties-diagonal}{source}{}{sourcepointer-4}{}}
\newlabel{lemma-properties-diagonal@cref}{{[originallemma][0][]lemma-properties-diagonal}{[1][0][]source}{}{}{}}
\makeatother
\crefname{originalreference}{original reference}{original references}
\Crefname{originalreference}{Original reference}{Original references}
\makeatletter
\newlabel{spaces-morphisms-lemma-characterize-proper}{{spaces-morphisms-lemma-characterize-proper}{source}{}{sourcepointer-5}{}}
\newlabel{spaces-morphisms-lemma-characterize-proper@cref}{{[originalreference][0][]spaces-morphisms-lemma-characterize-proper}{[1][0][]source}{}{}{}}
\makeatother
\title{Uniqueness in the valuative criterion implies that the diagonal is separated}
\author{The Stacks Project Authors}
\date{stacks-project:stacks-morphisms.tex:local-snapshot-20261012; GNU FDL 1.2 or later}
\begin{document}
\maketitle
\noindent\textit{Editorial note.} All delivered mathematical text is the author's own, copied from the pinned Stacks Project chapter source sources/stacks-algebraic-stacks-morphisms/pinned-original-source.tex; no statement or proof step was rewritten, reordered or removed.\par\medskip
\noindent\textit{Editorial note.} Preamble provenance: the Stacks Project keeps one single shared preamble for every chapter; this card takes it directly from the staged sources/stacks-algebraic-stacks-morphisms/preamble.tex as three verbatim spans (lines 8-12, 14-18 and 154-217), so no derived file was written and nothing was added or redefined. The wrapper class and the base packages are supplied by the assembler.\par\medskip
\noindent\textit{Editorial note.} Bibliography provenance: neither the counted statement nor its proof contains a citation (this lane rescans the whole span before writing the card), so the card carries no bibliography and no entry was reconstructed.\par\medskip
\noindent\textit{Editorial note.} Selection provenance: this card was opened under the file-based control contract (CONTROL.md), with the mechanical pre-card gate run on its statement and proof interval before the card was written.\par\medskip
\noindent\textit{Allowed standard prerequisites (editorial).}\begin{itemize}
\item The translation of uniqueness in the valuative criterion into a condition on the category of dotted arrows. Used throughout the counted proof (source lines 9633-9650). Turns the hypothesis into a statement about morpisms of dotted arrows.
\item The description of the diagonal as a morphism satisfying the corresponding valuative criterion. Invoked inside the counted proof (source lines 9651-9668). Reduces separatedness of the diagonal to the uniqueness property for the diagonal.
\item The valuative criterion for separatedness of a morphism of algebraic stacks. Invoked at the end of the counted proof (source lines 9669-9680). Concludes separatedness from the established uniqueness property.
\end{itemize}\medskip
\ifdefined\hypertarget\hypertarget{sourcepointer-1}{}\fi
\noindent\textit{Source reference.} Original definition definition-fill-in-diagram: sources/stacks-algebraic-stacks-morphisms/pinned-original-source.tex, source lines 9061--9061. Navigation pointer only: the counted text refers to the Stacks item definition-fill-in-diagram, which is an existing labeled result of this same chapter and is not part of this card. The pointer carries the target's own label as its number and its definition line as its span; its content is not reproduced here.\par\medskip
\ifdefined\hypertarget\hypertarget{sourcepointer-2}{}\fi
\noindent\textit{Source reference.} Original lemma lemma-cat-dotted-arrows: sources/stacks-algebraic-stacks-morphisms/pinned-original-source.tex, source lines 9115--9115. Navigation pointer only: the counted text refers to the Stacks item lemma-cat-dotted-arrows, which is an existing labeled result of this same chapter and is not part of this card. The pointer carries the target's own label as its number and its definition line as its span; its content is not reproduced here.\par\medskip
\ifdefined\hypertarget\hypertarget{sourcepointer-3}{}\fi
\noindent\textit{Source reference.} Original lemma lemma-existence-representable: sources/stacks-algebraic-stacks-morphisms/pinned-original-source.tex, source lines 9409--9409. Navigation pointer only: the counted text refers to the Stacks item lemma-existence-representable, which is an existing labeled result of this same chapter and is not part of this card. The pointer carries the target's own label as its number and its definition line as its span; its content is not reproduced here.\par\medskip
\ifdefined\hypertarget\hypertarget{sourcepointer-4}{}\fi
\noindent\textit{Source reference.} Original lemma lemma-properties-diagonal: sources/stacks-algebraic-stacks-morphisms/pinned-original-source.tex, source lines 344--344. Navigation pointer only: the counted text refers to the Stacks item lemma-properties-diagonal, which is an existing labeled result of this same chapter and is not part of this card. The pointer carries the target's own label as its number and its definition line as its span; its content is not reproduced here.\par\medskip
\ifdefined\hypertarget\hypertarget{sourcepointer-5}{}\fi
\noindent\textit{Source reference.} Original reference spaces-morphisms-lemma-characterize-proper: sources/stacks-algebraic-stacks-morphisms/pinned-original-source.tex, source lines 9176--9176. Navigation pointer only: the counted text refers to the Stacks item spaces-morphisms-lemma-characterize-proper, which belongs to another chapter of the project and is not part of this card. The pointer carries the target's own label as its number because this chapter's staged snapshot does not contain the other chapter's tag table; its content is not reproduced here.\par\medskip
\begin{lemma}
\label{lemma-converse-uniqueness-and-diagonal}
Let $f : \mathcal{X} \to \mathcal{Y}$ be a morphism of algebraic stacks.
If $f$ is separated, then $f$ satisfies the
uniqueness part of the valuative criterion.
\end{lemma}

\begin{proof}
Since $f$ is separated we see that all categories of dotted arrows
are setoids by Lemma \ref{lemma-cat-dotted-arrows}.
Consider a diagram
$$
\xymatrix{
\Spec(K) \ar[r]_-x \ar[d]_j & \mathcal{X} \ar[d]^f \\
\Spec(A) \ar[r]^-y \ar@{..>}[ru] & \mathcal{Y}
}
$$
and a $2$-morphism $\gamma : y \circ j \to f \circ x$ as in
Definition \ref{definition-fill-in-diagram}. Consider two
objects $(a, \alpha, \beta)$ and $(a', \beta', \alpha')$
of the category of dotted arrows. To finish the proof we
have to show these objects are isomorphic. The isomorphism
$$
f \circ a \xrightarrow{\beta^{-1}} y \xrightarrow{\beta'} f \circ a'
$$
means that $(a, a', \beta' \circ \beta^{-1})$ is a morphism
$\Spec(A) \to \mathcal{X} \times_\mathcal{Y} \mathcal{X}$.
On the other hand, $\alpha$ and $\alpha'$ define
a $2$-arrow
$$
(a, a', \beta' \circ \beta^{-1}) \circ j =
(a \circ j, a' \circ j,
(\beta' \star \text{id}_j) \circ (\beta \star \text{id}_j)^{-1})
\xrightarrow{(\alpha, \alpha')} (x, x, \text{id}) = \Delta_f \circ x
$$
Here we use that both $(a, \alpha, \beta)$ and $(a', \alpha', \beta')$
are dotted arrows with respect to $\gamma$.
We obtain a commutative diagram
$$
\xymatrix{
\Spec(K) \ar[d]_j \ar[rr]_x & & \mathcal{X} \ar[d]^{\Delta_f} \\
\Spec(A) \ar[rr]^{(a, a', \beta' \circ \beta^{-1})} & &
\mathcal{X} \times_\mathcal{Y} \mathcal{X}
}
$$
with $2$-commutativity witnessed by $(\alpha, \alpha')$. Now
$\Delta_f$ is representable by algebraic spaces
(Lemma \ref{lemma-properties-diagonal})
and proper as $f$ is separated. Hence by
Lemma \ref{lemma-existence-representable}
and the valuative criterion for properness for algebraic spaces
(Morphisms of Spaces, Lemma \ref{spaces-morphisms-lemma-characterize-proper})
we see that there exists a dotted arrow.
Unwinding the construction, we see that this means
$(a, \alpha, \beta)$ and $(a', \alpha', \beta')$
are isomorphic in the category of dotted arrows as desired.
\end{proof}

\end{document}
